\rSec2[expected.unexpected]{Class template unexpected}

\rSec3[expected.un.general]{General}

\begin{codeblock}
template<class E>
class @\libglobal{unexpected}@ {
  // [expected.un.general] para 2: ill-formed instantiations

public:
  constexpr unexpected(const unexpected&) = default;
  constexpr unexpected(unexpected&&) = default;

  template<class Err = E>
    requires(!is_same_v<remove_cvref_t<Err>, unexpected> &&
             !is_same_v<remove_cvref_t<Err>, in_place_t> && is_constructible_v<E, Err>)
  constexpr explicit unexpected(Err&& e) noexcept(is_nothrow_constructible_v<E, Err>);

  template<class... Args>
    requires is_constructible_v<E, Args...>
  constexpr explicit unexpected(in_place_t, Args&&... args) noexcept(
      is_nothrow_constructible_v<E, Args...>);

  template<class U, class... Args>
    requires is_constructible_v<E, initializer_list<U>&, Args...>
  constexpr explicit unexpected(
      in_place_t, initializer_list<U> il,
      Args&&... args) noexcept(is_nothrow_constructible_v<E, initializer_list<U>&,
                                                          Args...>);

  constexpr unexpected& operator=(const unexpected&) = default;
  constexpr unexpected& operator=(unexpected&&) = default;

  constexpr const E& error() const& noexcept;
  constexpr E& error() & noexcept;
  constexpr const E&& error() const&& noexcept;
  constexpr E&& error() && noexcept;

  constexpr void swap(unexpected& other) noexcept(is_nothrow_swappable_v<E>);

  template<class E2>
  friend constexpr bool operator==(const unexpected& x, const unexpected<E2>& y);

  friend constexpr void swap(unexpected& x, unexpected& y) noexcept(noexcept(x.swap(y)))
    requires is_swappable_v<E>;

private:
  E @\exposidnc{unex}@; // exposition only
};
\end{codeblock}

\begin{itemdescr}
\pnum
\mandates
A program that instantiates the definition of \tcode{unexpected} for a non-object type other than an lvalue reference type, an array type, a specialization of \tcode{unexpected}, or a cv-qualified type is ill-formed.

\pnum
\remarks
Subclause \iref{expected.unexpected} describes the class template \tcode{unexpected} that represents unexpected objects stored in \tcode{expected} objects.
\end{itemdescr}

\rSec3[expected.un.cons]{Constructors}

\indexlibraryctor{unexpected}%
\begin{itemdecl}
template<class Err = E>
  requires(!is_same_v<remove_cvref_t<Err>, unexpected> &&
           !is_same_v<remove_cvref_t<Err>, in_place_t> && is_constructible_v<E, Err>)
constexpr explicit unexpected(Err&& e) noexcept(is_nothrow_constructible_v<E, Err>);
\end{itemdecl}

\begin{itemdescr}
\pnum
\constraints
\begin{itemize}
\item \tcode{is_same_v<remove_cvref_t<Err>, unexpected>} is \tcode{false}; and
\item \tcode{is_same_v<remove_cvref_t<Err>, in_place_t>} is \tcode{false}; and
\item \tcode{is_constructible_v<E, Err>} is \tcode{true}.
\end{itemize}

\pnum
\effects
Direct-non-list-initializes \tcode{unex} with \tcode{std::forward<Err>(e)}.

\pnum
\throws
Any exception thrown by the initialization of \tcode{unex}.
\end{itemdescr}

\indexlibraryctor{unexpected}%
\begin{itemdecl}
template<class... Args>
  requires is_constructible_v<E, Args...>
constexpr explicit unexpected(in_place_t, Args&&... args) noexcept(
    is_nothrow_constructible_v<E, Args...>);
\end{itemdecl}

\begin{itemdescr}
\pnum
\constraints
\tcode{is_constructible_v<E, Args...>} is \tcode{true}.

\pnum
\effects
Direct-non-list-initializes \tcode{unex} with \tcode{std::forward<Args>(args)...}.

\pnum
\throws
Any exception thrown by the initialization of \tcode{unex}.
\end{itemdescr}

\indexlibraryctor{unexpected}%
\begin{itemdecl}
template<class U, class... Args>
  requires is_constructible_v<E, initializer_list<U>&, Args...>
constexpr explicit unexpected(
    in_place_t, initializer_list<U> il,
    Args&&... args) noexcept(is_nothrow_constructible_v<E, initializer_list<U>&,
                                                        Args...>);
\end{itemdecl}

\begin{itemdescr}
\pnum
\constraints
\tcode{is_constructible_v<E, initializer_list<U>&, Args...>} is \tcode{true}.

\pnum
\effects
Direct-non-list-initializes \tcode{unex} with \tcode{il, std::forward<Args>(args)...}.

\pnum
\throws
Any exception thrown by the initialization of \tcode{unex}.
\end{itemdescr}

\rSec3[expected.un.obs]{Observers}

\indexlibrarymember{error}{unexpected}%
\begin{itemdecl}
constexpr const E& error() const& noexcept;
\end{itemdecl}

\begin{itemdescr}
\pnum
\returns
\tcode{unex}.
\end{itemdescr}

\indexlibrarymember{error}{unexpected}%
\begin{itemdecl}
constexpr E& error() & noexcept;
\end{itemdecl}

\indexlibrarymember{error}{unexpected}%
\begin{itemdecl}
constexpr const E&& error() const&& noexcept;
\end{itemdecl}

\begin{itemdescr}
\pnum
\returns
\tcode{std::move(unex)}.
\end{itemdescr}

\indexlibrarymember{error}{unexpected}%
\begin{itemdecl}
constexpr E&& error() && noexcept;
\end{itemdecl}

\rSec3[expected.un.swap]{Swap}

\indexlibrarymember{swap}{unexpected}%
\begin{itemdecl}
constexpr void swap(unexpected& other) noexcept(is_nothrow_swappable_v<E>);
\end{itemdecl}

\begin{itemdescr}
\pnum
\mandates
\tcode{is_swappable_v<E>} is \tcode{true}.

\pnum
\effects
Equivalent to: \tcode{using std::swap; swap(unex, other.unex);}
\end{itemdescr}

\indexlibraryglobal{swap}%
\begin{itemdecl}
friend constexpr void swap(unexpected& x, unexpected& y) noexcept(noexcept(x.swap(y)));
\end{itemdecl}

\begin{itemdescr}
\pnum
\constraints
\tcode{is_swappable_v<E>} is \tcode{true}.

\pnum
\effects
Equivalent to \tcode{x.swap(y)}.
\end{itemdescr}

\rSec3[expected.un.eq]{Equality operator}

\indexlibraryglobal{operator==}%
\begin{itemdecl}
template<class E2>
friend constexpr bool operator==(const unexpected& x, const unexpected<E2>& y);
\end{itemdecl}

\begin{itemdescr}
\pnum
\mandates
The expression \tcode{x.error() == y.error()} is well-formed and its result is convertible to \tcode{bool}.

\pnum
\returns
\tcode{x.error() == y.error()}.
\end{itemdescr}

\rSec3[expected.un.ref]{Partial specialization unexpected<E&>}

\begin{codeblock}
template<class E>
class unexpected<E&> {
public:
  constexpr unexpected(const unexpected&) = default;
  constexpr unexpected(unexpected&&) = default;
  template<class G = E>
    constexpr explicit unexpected(G&&) noexcept;
  template<class G = E>
    constexpr explicit unexpected(in_place_t, G&&) noexcept;

  constexpr unexpected& operator=(const unexpected&) = default;
  constexpr unexpected& operator=(unexpected&&) = default;

  constexpr E& error() const noexcept;

  constexpr void swap(unexpected& other) noexcept;

  template<class E2>
    friend constexpr bool operator==(const unexpected&, const unexpected<E2>&);

  friend constexpr void swap(unexpected& x, unexpected& y) noexcept;

private:
  E* unex;              // exposition only
};
\end{codeblock}

\begin{itemdescr}
\pnum
\mandates
A program that instantiates the definition of \tcode{unexpected<E&>} for an array type or a specialization of \tcode{unexpected} is ill-formed.

\pnum
\remarks
An object of type \tcode{unexpected<E&>} holds a pointer to an object of type \tcode{E}. The referenced object is not owned by the \tcode{unexpected<E&>} object. Unlike the primary template, \tcode{E} may be a cv-qualified type.
\end{itemdescr}

\indexlibraryctor{unexpected}%
\begin{itemdecl}
template<class G = E>
  requires(!is_same_v<remove_cvref_t<G>, unexpected> &&
           !is_same_v<remove_cvref_t<G>, in_place_t> && is_constructible_v<E&, G&&> &&
           !reference_constructs_from_temporary_v<E&, G>)
constexpr explicit unexpected(G&& e) noexcept;
\end{itemdecl}

\begin{itemdescr}
\pnum
\constraints
\begin{itemize}
\item \tcode{is_same_v<remove_cvref_t<G>, unexpected>} is \tcode{false},
\item \tcode{is_same_v<remove_cvref_t<G>, in_place_t>} is \tcode{false},
\item \tcode{is_constructible_v<E&, G>} is \tcode{true}, and
\item \tcode{reference_constructs_from_temporary_v<E&, G>} is \tcode{false}.
\end{itemize}

\pnum
\effects
Initializes \tcode{unex} with \tcode{addressof(static_cast<E&>(std::forward<G>(e)))}.

\pnum
\remarks
A constructor for which \tcode{reference_constructs_from_temporary_v<E&, G>} is \tcode{true} — one that would bind \tcode{E&} to a temporary — is defined as deleted.
\end{itemdescr}

\indexlibraryctor{unexpected}%
\begin{itemdecl}
template<class G = E>
  requires(is_constructible_v<E&, G&&> && !reference_constructs_from_temporary_v<E&, G>)
constexpr explicit unexpected(in_place_t, G&& e) noexcept;
\end{itemdecl}

\begin{itemdescr}
\pnum
\constraints
\tcode{is_constructible_v<E&, G>} is \tcode{true} and \tcode{reference_constructs_from_temporary_v<E&, G>} is \tcode{false}.

\pnum
\effects
Initializes \tcode{unex} with \tcode{addressof(static_cast<E&>(std::forward<G>(e)))}.
\end{itemdescr}

\indexlibrarymember{error}{unexpected}%
\begin{itemdecl}
constexpr E& error() const noexcept;
\end{itemdecl}

\begin{itemdescr}
\pnum
\returns
\tcode{*unex}.

\pnum
\remarks
The reference returned is not \tcode{const}-qualified even when \tcode{*this} is \tcode{const}; the constness of the \tcode{unexpected<E&>} object does not propagate to the referenced object.
\end{itemdescr}

\indexlibrarymember{swap}{unexpected}%
\begin{itemdecl}
constexpr void swap(unexpected& other) noexcept;
\end{itemdecl}

\begin{itemdescr}
\pnum
\effects
Exchanges \tcode{unex} and \tcode{other.unex}.
\end{itemdescr}

\indexlibraryglobal{operator==}%
\begin{itemdecl}
template<class E2>
friend constexpr bool operator==(const unexpected& x, const unexpected<E2>& y);
\end{itemdecl}

\begin{itemdescr}
\pnum
\mandates
The expression \tcode{x.error() == y.error()} is well-formed and its result is convertible to \tcode{bool}.

\pnum
\returns
\tcode{x.error() == y.error()}.
\end{itemdescr}

\indexlibraryglobal{swap}%
\begin{itemdecl}
friend constexpr void swap(unexpected& x, unexpected& y) noexcept;
\end{itemdecl}

\begin{itemdescr}
\pnum
\effects
Equivalent to \tcode{x.swap(y)}.
\end{itemdescr}
